% !TEX program = lualatex
\documentclass[11pt]{article}

\usepackage[a4paper,margin=2.4cm]{geometry}
\usepackage{amsmath,amssymb,amsfonts}
\usepackage{mathtools}
\usepackage{graphicx}
\usepackage{float}
\usepackage{booktabs}
\usepackage{array}
\usepackage{caption}
\usepackage{subcaption}
\usepackage{titlesec}
\usepackage{enumitem}
\usepackage[hidelinks]{hyperref}
\usepackage{xcolor}

% --- fonts: body Latin Modern; CJK fallback via Noto Sans TC (for stray 中文) ---
\usepackage{fontspec}
\setmainfont{Latin Modern Roman}
\setsansfont{Latin Modern Sans}
\setmonofont{Latin Modern Mono}
\usepackage{luatexja-fontspec}
\setmainjfont{Noto Sans TC}
\setsansjfont{Noto Sans TC}

% --- figures live in the two simulation folders ---
\graphicspath{%
  {../part1_glucose/figures/}%
  {../part2_hiv/figures/}%
}

\setlength{\parindent}{0pt}
\setlength{\parskip}{0.35em}
\titlespacing*{\section}{0pt}{1.1em}{0.4em}
\titlespacing*{\subsection}{0pt}{0.7em}{0.25em}
\captionsetup{font=small,labelfont=bf,skip=3pt}

% --- tighten inter-float spacing to lift layout density ---
\setlength{\floatsep}{6pt plus 2pt minus 2pt}
\setlength{\textfloatsep}{8pt plus 2pt minus 2pt}
\setlength{\intextsep}{6pt plus 2pt minus 2pt}
\setlength{\abovecaptionskip}{3pt}
\setlength{\belowcaptionskip}{2pt}

% shorthands
\newcommand{\degC}{^\circ\mathrm{C}}
\newcommand{\mgdl}{\ensuremath{\,\mathrm{mg}/100\,\mathrm{ml}}}
\newcommand{\uU}{\ensuremath{\,\mu\mathrm{U}/\mathrm{ml}}}
\newcommand{\Rzero}{\mathcal{R}_0}
\newcommand{\Reff}{\mathcal{R}_{\mathrm{eff}}}
\newcommand{\dd}[2]{\frac{\mathrm{d}#1}{\mathrm{d}#2}}

% lightweight shims so we don't need the (uninstalled) siunitx package
\newcommand{\si}[1]{\ensuremath{#1}}
\providecommand{\micro}{\mu}

\title{\vspace{-1.2cm}\textbf{BME5113 Biological Systems Modeling and Analysis}\\[4pt]
\large Term Project --- Hormonal Glucose Control and HIV Vaccination Modeling}
\author{Chuan-Han Li (李傳漢)\quad B11611027\\ \small Department of Biomechatronics Engineering, National Taiwan University}
\date{June 2026}

\begin{document}
\maketitle
\thispagestyle{empty}

\begin{abstract}
\noindent This report addresses the two parts of the term project. \textbf{Part~1} implements the
Cobelli \emph{et~al.}\ (1982) intermediate-complexity glucose--insulin--glucagon model (Haefner
Ch.~12) and works through exercises 1--7: reproducing the IVGTT response, tracing where ingested
glucose goes, explaining the insulin ``hump'', and simulating diabetic, obese, hyperinsulinaemic and
insulin-overdose subjects. A central practical finding is that the literal Table~12.2 parameters do
not reproduce Fig.~12.4 and require a small, explicitly documented calibration. \textbf{Part~2}
implements the simplified Imperial College (``sIC'') HIV/AIDS compartment model (Haefner Ch.~15),
extends it with a waning vaccine (protected compartments $P_{f,2},P_{m,2}$, vaccination rate
$\nu=0.65\,$yr$^{-1}$, loss of protection $l=0.1\,$yr$^{-1}$), and studies the epidemic peak/decline,
cost, optimum vaccination rate, and vaccination versus condom use. Two modeling judgments are
required and documented: the literal progression rate $\gamma=1.16\,$yr$^{-1}$ gives $\Rzero<1$ and
must be corrected to a biologically grounded $\gamma=0.1\,$yr$^{-1}$ ($\Rzero\approx2.35$); and
protected individuals must remain in the frequency-dependent partner pool for the vaccine to express
herd immunity. All simulations are in Python (\texttt{scipy.integrate.solve\_ivp}); key results are
cross-checked with a second integrator.
\end{abstract}

\tableofcontents
\vspace{0.6em}

\section{Part 1 --- Hormonal Control of Glucose}

This part implements the Cobelli \emph{et~al.}\ (1982) intermediate-complexity model of
glucose--insulin--glucagon regulation, as presented in Haefner Ch.~12, and works through
exercises 1--7 of \S12.5. The core is written in Python; integration uses
\texttt{scipy.integrate.solve\_ivp} (LSODA), with a Radau cross-check. All exercises run against a
single verified core (\texttt{cobelli.py}). Before the exercises I describe the model and the one
non-trivial step needed to make it run at all: a small, documented calibration of an under-specified
parameter table.

\subsection{The Cobelli intermediate-complexity model}

The model tracks seven state variables in three coupled subsystems, structured as a Forrester
negative-feedback diagram (Fig.~12.2 of the text):

\begin{itemize}[leftmargin=1.4em,itemsep=1pt]
\item \textbf{Glucose subsystem:} $g$ = glucose in plasma and extracellular fluid (mg).
\item \textbf{Glucagon subsystem:} $c$ = glucagon in plasma and interstitial fluid (nU).
\item \textbf{Insulin subsystem (five compartments):} $s$ = stored pancreatic insulin, $r$ =
promptly releasable pancreatic insulin, $l$ = liver insulin, $p$ = plasma insulin, $i$ =
interstitial-fluid insulin (all \si{\micro U}).
\end{itemize}

Each state variable has an associated concentration $\bar{g}=g/V_b$, $\bar{p}=p/V_p$, and so on,
where the volumes $V_b,V_p,V_l,V_i$ are fixed fractions of body weight. Feedback is expressed through
fluxes built from the hyperbolic tangent: a stimulatory effect uses
$0.5\,[1+\tanh(b(\Delta\bar{x}+c))]$ and an inhibitory effect uses $0.5\,[1-\tanh(b(\Delta\bar{x}+c))]$,
where $\Delta\bar{x}$ is the deviation of a concentration from its basal value, $b$ sets the slope and
$c$ the centre. Because $\tanh$ has range $\pm1$, each gate is bounded in $[0,1]$ and saturates
smoothly; multiplying several gates implements ``multiple limiting factors''. The governing equations
(Eqs.~12.1--12.7) are

\begin{align}
\dd{g}{t} &= \mathrm{NHGB} - F_3 - F_4 - F_5 + I_g(t), \\
\dd{c}{t} &= -h_{02}\,c + F_7, \\
\dd{i}{t} &= -m_{13}\,i + m_{31}\,p, \\
\dd{l}{t} &= -(m_{02}+m_{12})\,l + m_{21}\,p + F_6, \\
\dd{p}{t} &= -(m_{01}+m_{21}+m_{31})\,p + m_{12}\,l + m_{13}\,i + I_p(t), \\
\dd{r}{t} &= k_{21}\,s - k_{12}\,r - F_6, \\
\dd{s}{t} &= -k_{21}\,s + k_{12}\,r + W,
\end{align}

where $\mathrm{NHGB}=F_1-F_2$ is the net hepatic glucose balance (liver production minus uptake),
$F_3$ is renal excretion, $F_4$ peripheral (muscle/adipose) use, $F_5$ the nearly
insulin-independent CNS/red-blood-cell sink, $F_6$ pancreatic insulin secretion, $F_7$ glucagon
secretion, $W$ insulin synthesis into the stored pool, and $I_g,I_p$ are the glucose/insulin input
functions used to drive diagnostic tests. The secretion flux
$F_6=0.5\,a_6[1+\tanh(b_6(\Delta\bar{g}+c_6))]\,r$ is proportional to the releasable pool $r$, which
matters for the ``hump'' in Exercise~2.

\textbf{Transcription fix.} The renal-excretion gate is printed in the book as
$M_{31}=0.5[1+\tanh(b_{13}(\bar{g}+c_{13}))]$, but $b_{13}=2.75\times10^{-2}$, $c_{13}=20$ are glucose
self-suppression constants for $M_1$. Table~12.2 instead lists $b_{31}=20$, $c_{31}=-180$, which set a
renal threshold near $180\mgdl$ --- the physiologically correct reading. We therefore use
$b_{31},c_{31}$ (so $F_3$ stays near zero until glucose exceeds the threshold), treating the printed
$b_{13}$ as a typo.

\subsection{Calibration: reproducing Fig.~12.4 from an under-specified table}
\label{sec:calibration}

Taken literally and in the amount-based form of Eqs.~12.1--12.7, the Table~12.2 parameters do
\emph{not} reproduce Figs.~12.4a/b. This was verified independently and is a genuine gap in the
printed table, not a units error on our side (an amount$\leftrightarrow$concentration $V_b\approx140$
factor closes less than 1\% of it). Quantitatively, for a 70\,kg subject:

\begin{itemize}[leftmargin=1.4em,itemsep=1pt]
\item The literal basal glucose clearance is $F_3+F_4+F_5\approx8.4\times10^{-4}\,$mg/min against a
basal glucose amount of $12{,}810$\,mg, giving a clearance timescale of $\sim1.5\times10^{7}$\,min.
Glucose essentially never clears (Fig.~12.4 shows recovery in $\sim$90\,min).
\item The basal plasma insulin amount is $34{,}650\,\si{\micro U}$, which needs a basal secretion of
$F_6\approx1.6\times10^{4}\,\si{\micro U}$/min, but the literal synthesis flux is
$W(\text{basal})\approx0.017\,\si{\micro U}$/min --- about $10^{6}$ times too small. The insulin and
glucagon axes collapse to zero and no IVGTT response is possible.
\end{itemize}

The model is explicitly described as ``semi-phenomenological\,\dots\ tailored/scaled to a particular
patient'' (Haefner \S12.3). Sitting the model at a patient's basal operating point by choosing the
under-determined source/clearance scales is exactly its intended use, so we applied a minimal,
structure-preserving calibration and changed \emph{no} feedback shape:

\begin{enumerate}[leftmargin=1.5em,itemsep=2pt]
\item A uniform \texttt{glucose\_flux\_scale} (default $10^{5}$) multiplies the four glucose fluxes
(NHGB, $F_3$, $F_4$, $F_5$), setting the clearance timescale to $\sim$90\,min. Because
$\mathrm{NHGB}=F_3+F_4+F_5$ at basal, a \emph{common} scale cancels, so the basal steady state and
every $\tanh$ feedback shape are exactly unchanged (the basal residual stays at $3.6\times10^{-12}$
for scale $\in\{1,10^{3},10^{5},10^{7}\}$; setting it to $1.0$ recovers the literal table).
\item The source scales $a_w$ (insulin synthesis) and $a_{71}$ (glucagon secretion) are chosen so the
published clinical basal is an exact steady state. No $\tanh$ slope $b$, centre $c$, or transfer rate
($m_{ij}$, $k_{ij}$) is altered.
\item $b_{52}$, a constant offset in the $F_5$ CNS/RBC sink, is shifted by the tiny published
NHGB$-(F_3+F_4+F_5)$ imbalance so that $\dd{g}{t}=0$ at basal exactly. This is a constant offset,
not a shape change.
\end{enumerate}

After calibration the basal residual is $3.6\times10^{-12}$; a normal IVGTT gives a glucose peak of
$250\mgdl$ recovering in $\sim$52\,min, plasma insulin rising $11\to40\uU$ then decaying, and glucagon
dipping then recovering, with LSODA and Radau agreeing to $1.6\times10^{-4}$. These match Fig.~12.4
qualitatively (Fig.~\ref{fig:verify}).

\textbf{What is not claimed.} We do not claim that $10^{5}$, or the recalibrated $a_w,a_{71},b_{52}$,
are the original Cobelli (1982) coefficients. They are a documented fit that reproduces the published
figures from an under-specified table while leaving all feedback structure intact.

\begin{figure}[t]
\centering
\includegraphics[width=0.62\linewidth]{verify_normal_ivgtt.png}
\caption{Verification of the calibrated core against a normal IVGTT. Plasma glucose peaks near
$250\mgdl$ and recovers within $\sim$52\,min; plasma insulin rises and decays; LSODA and Radau agree
to $1.6\times10^{-4}$. The calibration preserves the basal steady state and all $\tanh$ feedback
shapes.}
\label{fig:verify}
\end{figure}

\subsection{Exercise 1: Normal IVGTT and where glucose goes}

\textbf{Question.} Code the model, reproduce Figs.~12.4a/b (plasma glucose; plasma/liver/interstitial
insulin), additionally plot glucagon and the liver uptake $F_2$, and answer: where does the majority
of glucose go?

\textbf{Approach.} A normal 70\,kg subject receives a single IV glucose bolus of $0.33\,$g/kg
($=23.1$\,g) over 2\,min, integrated from the calibrated basal state for 150\,min. The glucose-disposal
fluxes are integrated over the window to apportion the load.

Plasma glucose rises from basal $91.5$ to a peak of $250\mgdl$ at $t=2$\,min, recovering within 5\% of
basal by $t=52$\,min. Plasma insulin rises from $11$ to a peak of $40\uU$ within minutes, and glucagon
is suppressed from $75$ down to $6$\,nU/ml by the combined glucose and insulin rise
(Figs.~\ref{fig:ex1_glucose}--\ref{fig:ex1_glucagon}). The cumulative disposal split is given in
Table~\ref{tab:ex1_split}.

\begin{figure}[H]
\centering
\begin{subfigure}{0.48\linewidth}
\includegraphics[width=\linewidth]{ex1_glucose.png}
\caption{Plasma glucose (Fig.~12.4a).}
\end{subfigure}\hfill
\begin{subfigure}{0.48\linewidth}
\includegraphics[width=\linewidth]{ex1_insulin.png}
\caption{Plasma / liver / interstitial insulin (Fig.~12.4b).}
\end{subfigure}
\caption{Normal IVGTT ($23.1$\,g over 2\,min). Glucose peaks at $250\mgdl$ then recovers; all three
insulin compartments rise within minutes of the bolus.}
\label{fig:ex1_glucose}
\end{figure}

\begin{table}[H]
\centering
\caption{Cumulative glucose disposal over the 150-min IVGTT response. Peripheral disposal
$F_3+F_4+F_5=40{,}170$\,mg.}
\label{tab:ex1_split}
\begin{tabular}{lrr}
\toprule
Sink & Cumulative (mg) & \% of peripheral disposal \\
\midrule
$F_5$ CNS / RBC & $24{,}503$ & $61.0$ \\
$F_4$ muscle / adipose & $11{,}202$ & $27.9$ \\
$F_3$ renal & $4{,}465$ & $11.1$ \\
\midrule
$\int\mathrm{NHGB}$ (signed) & $16{,}855$ & --- \\
\bottomrule
\end{tabular}
\end{table}

\textbf{Discussion.} The majority of disposed glucose goes to the CNS and red blood cells ($F_5$,
$61.0\%$), followed by insulin-sensitive muscle and adipose tissue ($F_4$, $27.9\%$); renal excretion
($F_3$) contributes $11.1\%$. This is physiologically sensible: $F_5$ is a large, nearly
insulin-independent sink, while $F_3$ stays near zero until $\bar{g}$ exceeds the $\sim$180\,mg/100\,ml
renal threshold --- here the peak is $250\mgdl$, so $F_3$ is engaged only modestly. In Forrester terms,
the bolus drives glucose up, which opens $F_6$ ($\propto r$) so insulin climbs, which both clears
glucose peripherally and (via inhibitory $-\tanh$ gates) suppresses glucagon. The sign of
$\int\mathrm{NHGB}$ ($+16{,}855$\,mg) indicates the liver is a net source over the full window, as the
basal hepatic output dominates once glucose has cleared.

\begin{figure}[H]
\centering
\begin{subfigure}{0.48\linewidth}
\includegraphics[width=\linewidth]{ex1_glucagon_F2.png}
\caption{Glucagon and liver uptake $F_2$.}
\end{subfigure}\hfill
\begin{subfigure}{0.48\linewidth}
\includegraphics[width=\linewidth]{ex1_disposal_split.png}
\caption{Cumulative disposal fluxes.}
\end{subfigure}
\caption{Glucagon is suppressed during the excursion; the cumulative-flux panel shows CNS/RBC ($F_5$)
as the dominant sink.}
\label{fig:ex1_glucagon}
\end{figure}

\subsection{Exercise 2: The insulin ``hump''}

\textbf{Question.} Why does the ``hump'' in insulin concentration develop in Fig.~12.5 (repeated
glucose pulses, normal subject), and why is there a similar hump for obese persons in Fig.~12.7?

\textbf{Approach.} (i) A normal subject is given an escalating repeated-pulse train
($[0.5,1.0,2.5,5.0,10.0,20.0,40.0]$\,g at $t=[0,30,70,110,180,270,400]$\,min) over 700\,min. (ii) An
obese subject is given a single $23.1$\,g IVGTT from its pre-equilibrated basal.

In the normal repeated-pulse run a hump is detected: insulin re-rises to $\sim$37\uU\ near $t=450$\,min
after dipping to $\sim$31\uU\ (Fig.~\ref{fig:ex2}). In the obese single-IVGTT run a hump is likewise
detected, with a secondary insulin rise to $\sim$32\uU\ near $t=105$\,min (dip $\sim$22\uU; peak
$36\uU$ vs basal $11\uU$).

\textbf{Mechanism.} Both humps share one origin: secretion
$F_6=0.5\,a_6[1+\tanh(b_6(\Delta\bar{g}+c_6))]\,r$ is proportional to the releasable pool $r$, which
refills from the stored pool $s$ on its own slow timescale ($\dd{r}{t}=k_{21}s-k_{12}r-F_6$;
$\dd{s}{t}=-k_{21}s+k_{12}r+W$). You cannot secrete insulin you do not have, and $r$ cannot refill
instantly. In Fig.~12.5 each pulse drains $r$; because pulses are spaced shorter than the
refill/clearance time, $r$ is repeatedly depleted and over-refilled from $s$, and the multi-compartment
lag ($r\to l\to p\to i$) stacks the delayed contributions, so plasma insulin does not decay
exponentially between pulses but accumulates into a delayed hump after the larger pulses. In the obese
case (Fig.~12.7) the same pool dynamics are driven by parameters rather than repeated input: $b_6=0.5$
(vs nominal $9.23\times10^{-2}$) steepens the $F_6$ gate so a single load triggers an outsized
drain-and-refill of $r$, and $m_{02}=0.13$ slows downstream clearance, so the pools overshoot and
secretion re-engages as $r$ refills --- an abnormal hump on the insulin-decay curve even though the
glucose response is near-normal.

\begin{figure}[H]
\centering
\begin{subfigure}{0.48\linewidth}
\includegraphics[width=\linewidth]{ex2_normal_hump.png}
\caption{Normal, repeated pulses.}
\end{subfigure}\hfill
\begin{subfigure}{0.48\linewidth}
\includegraphics[width=\linewidth]{ex2_obese_hump.png}
\caption{Obese, single IVGTT.}
\end{subfigure}
\caption{The insulin hump in both cases. Panels show glucose, insulin, the $r/s$ pools and $F_6$: the
depletion--refill cycling of $r$ drives the secondary insulin rise.}
\label{fig:ex2}
\end{figure}

\subsection{Exercise 3: Repeated pulses, diabetic vs normal}

\textbf{Question.} Apply the Fig.~12.5 pulse sequence to a diabetic subject and compare with a normal
subject.

\textbf{Approach.} The identical escalating train is applied to both over 800\,min using the single
canonical core. The diabetic uses the Table~12.3 secretion/uptake-shape changes, with the
insulin-driven hepatic ($F_2$) and peripheral ($H_4$) uptake gates frozen, and a calibration to its
own elevated operating point.

\begin{table}[H]
\centering
\caption{Repeated-pulse response, normal vs diabetic.}
\label{tab:ex3}
\begin{tabular}{lrr}
\toprule
Quantity & Normal & Diabetic \\
\midrule
Basal glucose ($\mgdl$) & $91.5$ & $140.0$ \\
Peak glucose ($\mgdl$) & $366$ & $418$ \\
Final glucose ($\mgdl$) & $89.9$ & $139.7$ \\
Basal insulin ($\uU$) & $11.0$ & $7.0$ \\
Peak insulin ($\uU$) & $62$ & $18.3$ \\
Recovery after last pulse (min) & $462$ & $524$ \\
\bottomrule
\end{tabular}
\end{table}

\textbf{Disclosed modeling choice.} Table~12.3 lists $F_2=0.037$, but in the textbook's native
(un-scaled) units. Fed through our $\texttt{glucose\_flux\_scale}=10^{5}$, the literal value becomes a
$\sim$3700\,mg/min hepatic ``vacuum'' ($\sim$70$\times$ the normal basal $F_2\approx52$\,mg/min) that
clamps glucose flat --- physiologically backwards. The scale-consistent reading of ``$F_2/H_2$ replaced
by a constant'' is that the diabetic liver can no longer \emph{ramp} uptake with insulin, so we freeze
$F_2$ near the normal basal hepatic-uptake flux and freeze $H_4$ at the Table~12.3 value $0.0012$. The
elevated basal does not emerge from freezing alone (the $M_1$ glucose self-suppression and $G_1$
glucagon gates make $91.5\mgdl$ a strong attractor), so the diabetic is calibrated to its own operating
point ($\mathrm{DIABETIC\_BASAL}=140/7/75$) using the same under-determined scales ($a_w,a_{71},b_{52}$)
as the normal subject --- no feedback shape is changed. At the low liver-insulin diabetic state the
\emph{unfrozen} hepatic-production gate $H_1$ naturally sits high ($\sim$0.66 vs $0.25$ normal):
hepatic production is no longer insulin-suppressed, the textbook diabetic mechanism, emergent rather
than forced.

\textbf{Discussion.} The diabetic basal glucose ($140\mgdl$) is elevated and its basal insulin
($7.0\uU$) blunted, both relative to normal ($91.5\mgdl$, $11.0\uU$). Its peak insulin ($18.3\uU$) is
far below the normal $62\uU$ --- the blunted Table~12.3 secretion machinery cannot mount the same
overshoot, so the characteristic insulin hump is suppressed. With impaired uptake and unsuppressed
production, the diabetic clears each pulse more slowly and stacks glucose above its already-high basal.
On a standard IVGTT this same configuration recovers in $\sim$100\,min vs $\sim$52\,min for normal
($\sim$1.9$\times$), the model expression of the textbook's ``$\sim$2$\times$ slower recovery''
(\S12.3.6).

\begin{figure}[t]
\centering
\includegraphics[width=0.66\linewidth]{ex3_diabetic_vs_normal.png}
\caption{Repeated-pulse plasma glucose and insulin, normal vs diabetic (overlay and side-by-side). The
diabetic stays elevated above its higher basal with a blunted, hump-free insulin response.}
\label{fig:ex3}
\end{figure}

\subsection{Exercise 4: Meal (60-min infusion) vs IVGTT pulse}

\textbf{Question.} Deliver glucose not as an instantaneous IVGTT pulse but as a constant input spread
over 60\,min (a meal): 25\,g to a 70\,kg subject. Compare with an IVGTT of the same total dose.

\textbf{Approach.} A normal subject receives $25$\,g either as a constant infusion
$I_g=416.7$\,mg/min over $[0,60]$\,min, or as a $\sim$2-min bolus ($I_g=12{,}500$\,mg/min). The only
difference between runs is the temporal spread of the input. The 60-min rectangular infusion is a
deliberate first-order stand-in for the smooth, delayed gastrointestinal rate of appearance $R_a(t)$ of
a real meal (cf.\ Dalla Man \emph{et~al.}\ 2007); the full GI submodel is not required here.

\begin{table}[H]
\centering
\caption{Meal (60-min infusion) versus IVGTT (2-min bolus), same 25\,g total. Basal glucose
$91.5\mgdl$, insulin $11.0\uU$.}
\label{tab:ex4}
\begin{tabular}{lrr}
\toprule
Metric & Meal (60\,min) & IVGTT (2\,min) \\
\midrule
Peak plasma glucose ($\mgdl$) & $143.9$ & $263.4$ \\
Time-to-peak glucose (min) & $59.9$ & $2.0$ \\
Peak plasma insulin ($\uU$) & $30.5$ & $40.2$ \\
Time-to-peak insulin (min) & $13.8$ & $4.2$ \\
Glucose recovery (min) & $82.5$ & $53.9$ \\
\bottomrule
\end{tabular}
\end{table}

\textbf{Discussion.} The IVGTT dumps all 25\,g into the glucose pool within 2\,min, faster than any
clearance flux can respond, so glucose spikes to $263.4\mgdl$ almost immediately. The meal delivers the
identical mass at only $417$\,mg/min --- slow enough that the clearance fluxes ($F_3,F_4,F_5$) and the
suppression of hepatic production remove glucose as it arrives, giving a much lower, broader excursion
peaking at $143.9\mgdl$ near the end of the infusion. The insulin response mirrors the drive: the
gentler glucose rise opens the $W$ and $F_6$ secretion gates less strongly, so the meal insulin peak
($30.5\uU$) is blunted relative to the sharp IVGTT spike ($40.2\uU$), and then settles to a lower
sustained plateau as $r$ partially depletes and is refilled by synthesis $W$. This is why an oral meal
stresses glucose homeostasis far less than an equal-dose IV bolus, and why tolerance tests use a rapid
bolus to maximally probe the feedback system.

\begin{figure}[t]
\centering
\includegraphics[width=0.66\linewidth]{ex4_meal_vs_ivgtt.png}
\caption{Meal vs IVGTT: plasma glucose (overlay), plasma insulin (overlay), and the input rate
$I_g(t)$. Same 25\,g total, very different temporal spread.}
\label{fig:ex4}
\end{figure}

\subsection{Exercise 5: Hyperinsulinism and IV-glucose rescue}

\textbf{Question.} (a) Simulate hyperinsulinism (first guess: increase $a_6$) and terminate the patient
when plasma glucose falls below $20\mgdl$; report the fatal setting and time of death. (b) Resuscitate
with IV glucose: how much is needed?

\textbf{Approach (disclosed reformulation).} Simply raising $a_6$ does not kill the subject. At steady
state the pancreatic mass balance forces $F_6=W$, and the $W$ gate closes as glucose falls, so
endogenous insulin saturates ($\sim$40\uU) and glucose only dips to $\sim$78\mgdl\ before recovering.
This homeostasis is correct --- a normal pancreas cannot kill you --- so clinical insulin shock requires
an \emph{autonomous} source (insulinoma or overdose) that does not switch off when glucose drops. We
model this faithfully as an autonomous insulinoma with output
$I_{p,\mathrm{tumor}}=(M-1)\,F_{6,\mathrm{basal}}$ \si{\micro U}/min ($F_{6,\mathrm{basal}}=1.638\times10^{4}$
\si{\micro U}/min), where multiplier $M=1$ is the healthy subject. The native pancreas (the $W=F_6$
balance, all $\tanh$ shapes, all transfer rates) is unchanged; death is a terminal event on
$\bar{g}-20$.

\begin{table}[H]
\centering
\caption{Fatal-hyperinsulinism sweep. The first fatal multiplier is $4.5\times$; the demo scenario
uses $8\times$.}
\label{tab:ex5}
\begin{tabular}{rrrcr}
\toprule
$a_6$ mult & $I_{p,\mathrm{tumor}}$ (\si{\micro U}/min) & min $\bar{g}$ ($\mgdl$) & died & $t_{\mathrm{death}}$ (min) \\
\midrule
$1.0$ & $0$ & $91.50$ & no & --- \\
$3.0$ & $3.277\times10^{4}$ & $79.66$ & no & --- \\
$4.0$ & $4.915\times10^{4}$ & $36.65$ & no & --- \\
$4.5$ & $5.735\times10^{4}$ & $20.10$ & yes & $79.6$ \\
$6.0$ & $8.192\times10^{4}$ & $20.22$ & yes & $49.6$ \\
$8.0$ & $1.147\times10^{5}$ & $20.11$ & yes & $43.3$ \\
$10.0$ & $1.475\times10^{5}$ & $20.24$ & yes & $38.9$ \\
\bottomrule
\end{tabular}
\end{table}

\textbf{Results.} The lethal threshold is the $4.5\times$ multiplier; the demonstration scenario
($8\times$) kills the subject at $t=43.3$\,min (Fig.~\ref{fig:ex5}). For the rescue, a continuous IV
dextrose drip is given over a defined 240-min clinical window and the total dose found by bisection
(keeping $\bar{g}\ge20$ throughout): the \textbf{minimum rescue dose is $69.11$\,g} of glucose over
240\,min. Below this the nadir still crosses 20 and the patient dies; at or above it the glucose floor
holds at 20 and the patient survives the window.

\textbf{Limitation.} Because the modelled insulinoma is permanent, IV glucose is temporizing, not
curative (clinically the tumour must be removed) --- exactly the textbook's ``short-term treatment =
large IV glucose''. The minimum dose is therefore quoted for a defined window; a longer window needs
proportionally more glucose. Quoting an unconditional ``grams to survive forever'' would be undefined
for a permanent source.

\begin{figure}[H]
\centering
\begin{subfigure}{0.48\linewidth}
\includegraphics[width=\linewidth]{ex5_fatal_decline.png}
\caption{Fatal decline ($8\times$).}
\end{subfigure}\hfill
\begin{subfigure}{0.48\linewidth}
\includegraphics[width=\linewidth]{ex5_rescue.png}
\caption{Un-rescued death vs rescued survival.}
\end{subfigure}
\caption{Hyperinsulinism. Left: glucose driven through the 70 and 20\,mg/100\,ml lines as insulin rises
far above basal. Right: the $69.11$\,g/240\,min drip holds the glucose floor at 20.}
\label{fig:ex5}
\end{figure}

\subsection{Exercise 6: Insulin bolus, normal vs obese}

\textbf{Question.} Simulate rapid ingestion of insulin at $0.10$\,U/kg over 2\,min, observe the
momentary hypoglycemia, and ask whether the subject dies; repeat for an obese subject.

\textbf{Approach.} The dose is $0.10\,\text{U/kg}\times70\,\text{kg}=7\,\text{U}=7.0\times10^{6}$
\si{\micro U}, delivered as a rectangular $I_p(t)$ pulse ($3.5\times10^{6}$ \si{\micro U}/min) into the
plasma-insulin compartment over 2\,min. Each subject starts from its own pre-equilibrated basal; death
is a terminal event at $\bar{g}<20\mgdl$.

\begin{table}[H]
\centering
\caption{Insulin-bolus response. Both subjects reach the same insulin peak; only the glucose nadir
differs.}
\label{tab:ex6}
\begin{tabular}{lrrrl}
\toprule
Subject & glucose nadir ($\mgdl$) & $t_{\mathrm{nadir}}$ (min) & insulin peak ($\uU$) & verdict \\
\midrule
Normal & $20.08$ & $21.3$ & $1541$ & DIES ($t=21.4$\,min) \\
Obese & $59.83$ & $37.2$ & $1542$ & survives ($t_{\mathrm{rec}}=103.6$\,min) \\
\bottomrule
\end{tabular}
\end{table}

\textbf{Discussion.} Both subjects receive the identical 7\,U bolus and reach essentially the same
plasma-insulin peak ($\sim$1541\uU), yet their glucose nadirs differ sharply ($20.1$ vs $59.8\mgdl$).
The difference comes entirely from peripheral disposal ($F_4$): the obese subject freezes the
interstitial-insulin gate $H_4$ at $0.0012$ (Table~12.3) and shifts the glucose gate ($b_{42},c_{42}$),
so the same insulin flood removes far less glucose per minute. This is the model's representation of
insulin resistance, and here it is \emph{protective}: the normal subject's fully insulin-responsive
tissue drives glucose all the way through the $20\mgdl$ line and dies, while the resistant obese subject
bottoms out in the hypoglycemic-but-survivable range and recovers as the bolus clears and
counter-regulation restores glucose. The normal result sits right at the lethal boundary, but the
terminal event fires consistently across solvers and step sizes, so the verdict is robust within the
model.

\begin{figure}[t]
\centering
\includegraphics[width=0.66\linewidth]{ex6_insulin_bolus_normal_vs_obese.png}
\caption{Insulin bolus, normal vs obese. Identical insulin peaks, but the normal nadir touches
$20\mgdl$ (death) while the insulin-resistant obese subject bottoms out at $\sim$60\mgdl\ and recovers.}
\label{fig:ex6}
\end{figure}

\subsection{Exercise 7: A bowl of vanilla ice cream}

\textbf{Question.} Simulate the glucose, insulin and glucagon dynamics for a normal, diabetic and obese
subject each consuming an average bowl of vanilla ice cream.

\textbf{Approach and assumption.} An average bowl is taken as $\sim$1 cup ($\sim$130\,g) of regular
vanilla ice cream. USDA FoodData Central lists $\sim$23--24\,g carbohydrate per 100\,g (almost entirely
sugar), giving $\sim$31\,g; we adopt \textbf{30\,g of digestible carbohydrate, treated as 100\%
appearing as plasma glucose} (a conservative upper bound, since lactose and fat slow real absorption).
The bowl is eaten over 40\,min, modelled as a constant infusion $I_g=750$\,mg/min on $[0,40]$\,min --- a
first-order stand-in for the gastrointestinal $R_a(t)$. The normal subject starts from its clinical
basal, the obese subject is settled to its (normal-anchored) basal, and the diabetic uses the same
canonical configuration as Exercise~3 (its own elevated basal $140/7/75$).

\begin{table}[H]
\centering
\caption{Ice-cream meal (30\,g over 40\,min): per-subject peaks and recovery. Recovery is the first
return within 5\% of that subject's own basal.}
\label{tab:ex7}
\begin{tabular}{lrrrrrr}
\toprule
Subject & basal $g$ & peak $g$ & basal ins & peak ins & min glucagon & $g$ recovery (min) \\
 & ($\mgdl$) & ($\mgdl$) & ($\uU$) & ($\uU$) & (nU/ml) & \\
\midrule
Normal & $91.5$ & $185.7$ & $11.0$ & $36.9$ & $2.2$ & $77.3$ \\
Diabetic & $140.0$ & $252.8$ & $7.0$ & $9.4$ & $1.5$ & $126.4$ \\
Obese & $91.5$ & $206.8$ & $11.0$ & $35.5$ & $1.2$ & $113.8$ \\
\bottomrule
\end{tabular}
\end{table}

\textbf{Discussion.} A single identical meal produces three distinct responses. The \emph{normal}
subject rises to $\sim$186\mgdl\ then clears back toward $91.5$ with a clean first-phase insulin spike
($11\to37\uU$). The \emph{obese} subject reaches a slightly higher glucose peak ($\sim$207\mgdl) with a
broadly similar excursion shape from the same $91.5$ basal --- ``nearly normal'' glucose, as the
textbook states --- but reproduces the characteristic abnormal hump in the insulin-decay curve (a
secondary rise after the first-phase peak, driven by the steep $F_6$ gate and faster $m_{02}$
clearance). The \emph{diabetic} subject starts from an elevated fasting basal ($140\mgdl$), rises to the
highest peak ($\sim$253\mgdl) and recovers most slowly ($\sim$126\,min vs $\sim$77\,min for normal),
with honestly blunted insulin ($7.0\to9.4\uU$): impaired insulin-driven uptake ($F_2,H_4$ frozen low)
plus unsuppressed hepatic production (high $H_1$) gives a higher, slower excursion. In all three
subjects the meal-driven glucose and insulin rise suppresses glucagon (the $F_7$ gates are negatively
driven by both), so glucagon plunges from $\sim$75 toward $\sim$1--2\,nU/ml during the meal and then
recovers.

\begin{figure}[t]
\centering
\begin{subfigure}{0.32\linewidth}
\includegraphics[width=\linewidth]{ex7_icecream_glucose.png}
\caption{Plasma glucose.}
\end{subfigure}\hfill
\begin{subfigure}{0.32\linewidth}
\includegraphics[width=\linewidth]{ex7_icecream_insulin.png}
\caption{Plasma insulin.}
\end{subfigure}\hfill
\begin{subfigure}{0.32\linewidth}
\includegraphics[width=\linewidth]{ex7_icecream_glucagon.png}
\caption{Plasma glucagon.}
\end{subfigure}
\caption{Ice-cream meal, three subjects. Glucose: diabetic highest and slowest. Insulin: normal clean
spike, obese abnormal hump, diabetic blunted. Glucagon: suppressed in all three during the meal.}
\label{fig:ex7}
\end{figure}

\subsection{Part 1 summary}

The Cobelli model reproduces the qualitative IVGTT behaviour of Fig.~12.4 once an under-specified
parameter table is calibrated to a patient's basal operating point --- a step that turned out to be the
main practical obstacle and is documented in full rather than hidden. With that single calibrated core,
exercises 1--7 recover the expected physiology: CNS/RBC tissue takes most of an ingested glucose load,
the insulin hump arises from finite-pool secretion ($F_6\propto r$), a spread-out meal stresses
homeostasis far less than an equal-dose IV bolus, and diabetic, obese, hyperinsulinaemic and
insulin-overdose subjects each show their characteristic signature. Two reformulations were needed and
disclosed honestly: hyperinsulinism requires an autonomous (insulinoma) source because the intact
pancreas is self-limiting, and the diabetic subject must be calibrated to its own elevated operating
point because the normal basal is a strong attractor. None of these changes the model's feedback
structure.

\clearpage
\section{Part 2 --- HIV Vaccination on the sIC AIDS Model}

This part takes the simplified Imperial College (sIC) AIDS compartment model from
Haefner~(2005), Chapter~15, and extends it with an HIV vaccination programme. The
task is essentially Exercise~15.6~\#5: vaccinate sexually active susceptibles into
a protected state at rate $\nu$, let protection wane at rate $l$, and then ask four
questions --- whether vaccination makes the epidemic peak and decline, what it
costs, what the optimum vaccination rate is, and how vaccination compares with
condom promotion. I work from the textbook equations and Table~15.2, the
vaccination design loosely follows Garnett \& Stover~(2002), and all numbers below
come from my own simulation of the model. Two of the modelling decisions needed to
make the textbook model behave as the textbook describes are not cosmetic, so I
state them explicitly and justify them rather than hide them.

\subsection{The sIC model}

The sIC model partitions the population into twelve compartments: the three disease
states susceptible ($S$), HIV-infected but pre-AIDS ($I$), and clinical AIDS ($A$),
crossed with two sexes (female $f$, male $m$) and two age classes (age~1 = birth to
15, pre-sexual; age~2 = 16 to death, sexually active). Only age-2 individuals
participate in transmission. The state vector is therefore
\[
\{S_{f,1},S_{f,2},S_{m,1},S_{m,2},\;
  I_{f,1},I_{f,2},I_{m,1},I_{m,2},\;
  A_{f,1},A_{f,2},A_{m,1},A_{m,2}\}.
\]
The model is a deliberate simplification of the full IC model: two sexes with
identical behaviour, a single AIDS stage, a single sexual-activity class (so the
social-mixing probability collapses to $\rho=1$), and two age classes.

\paragraph{Force of infection.}
HIV is sexually transmitted, so transmission is \emph{frequency-dependent}: people
seek a roughly fixed number of partners regardless of how large the population is,
so what matters is the \emph{proportion} of potential partners who are infectious,
not their absolute density. For a susceptible age-2 female the per-capita hazard is
(Eq.~15.6)
\begin{equation}
\lambda_{S_{f,2}} = c\,\rho\,\beta_{m,f}\,\frac{I_{m,2}}{S_{m,2}+I_{m,2}},
\qquad
\lambda_{S_{m,2}} = c\,\rho\,\beta_{f,m}\,\frac{I_{f,2}}{S_{f,2}+I_{f,2}},
\label{eq:foi}
\end{equation}
the second being the male analogue (Eq.~15.8). Here $c$ is the new-partner
acquisition rate, $\beta_{m,f}$ and $\beta_{f,m}$ are the per-partnership
transmission probabilities, and the factor $I/(S+I)$ is the frequency of
infectious individuals among possible partners. Two features deserve note. First,
the transmission probabilities are \emph{asymmetric}: $\beta_{m,f}=0.2$ versus
$\beta_{f,m}=0.075$, so male\,$\to$\,female transmission is about $2.7\times$ more
likely, a standard biological feature of heterosexual HIV. This is why women reach
higher equilibrium prevalence than men in the baseline run. Second, clinical-AIDS
individuals ($A$) are \emph{excluded} from the partner pool (assumed not sexually
active), so the denominator is $S+I$ only --- a point that turns out to matter for
the choice of epidemic seed.

\paragraph{The female equations.}
The full female subsystem (Eqs.~15.5a--f) is
\begin{align}
\dd{S_{f,1}}{t} &= \eta\theta\zeta\bigl(S_{f,2}+(1-\vartheta)I_{f,2}\bigr)
                   - \mu S_{f,1} - \xi S_{f,1}, \label{eq:Sf1}\\
\dd{S_{f,2}}{t} &= \xi S_{f,1} - (\lambda_{S_{f,2}}+\mu)\,S_{f,2}, \label{eq:Sf2}\\
\dd{I_{f,1}}{t} &= \eta\theta\zeta\vartheta\,I_{f,2}
                   - (\mu+\xi)I_{f,1} - \gamma I_{f,1}, \label{eq:If1}\\
\dd{I_{f,2}}{t} &= \lambda_{S_{f,2}} S_{f,2}
                   - (\mu+\gamma)I_{f,2} + \xi I_{f,1}, \label{eq:If2}\\
\dd{A_{f,1}}{t} &= \gamma I_{f,1} - (\mu+\xi+\alpha)A_{f,1}, \label{eq:Af1}\\
\dd{A_{f,2}}{t} &= \gamma I_{f,2} + \xi A_{f,1} - (\mu+\alpha)A_{f,2}. \label{eq:Af2}
\end{align}
The male equations~(15.7a--f) are identical in form, with the newborn sex split
$(1-\eta)$ replacing $\eta$ and the male force of infection
$\lambda_{S_{m,2}}$ replacing $\lambda_{S_{f,2}}$; note that male births are still
driven by the age-2 female compartments, since the model assumes males do not limit
reproduction.

The mechanisms encoded here are: \emph{reproduction}, with births created only from
age-2 females at fecundity $\theta$ (Eq.~\ref{eq:Sf1}); \emph{perinatal
transmission}, where an infected mother bears an infected newborn with probability
$\vartheta=0.35$ (feeding $I_{\cdot,1}$, Eq.~\ref{eq:If1}) and an uninfected
newborn with probability $1-\vartheta$ (feeding $S_{\cdot,1}$); \emph{ageing} at
rate $\xi$, moving each state from age~1 to age~2 within sex; \emph{AIDS
progression} at rate $\gamma$ moving $I\to A$; \emph{AIDS mortality} $\alpha$ added
to natural mortality $\mu$ in the $A$ compartments; and \emph{demography}, with
natural death $\mu$ acting on every compartment. The parameters are listed in
Table~\ref{tab:sic-params}.

\begin{table}[t]
\centering
\caption{Parameters of the simple IC (sIC) AIDS model, after Haefner~(2005),
Table~15.2 (values attributed there to Garnett \& Anderson, 1993). Variable
``values'' are the initial conditions used for the baseline run.}
\label{tab:sic-params}
\begin{tabular}{llrl}
\toprule
Symbol & Meaning & {Value} & Unit \\
\midrule
$\alpha$        & AIDS-stage extra death rate                 & 1.0    & yr$^{-1}$ \\
$\beta_{f,m}$   & female$\to$male transmission probability    & 0.075  & --- \\
$\beta_{m,f}$   & male$\to$female transmission probability     & 0.2    & --- \\
$c$             & new-partner acquisition rate                 & 2.35   & yr$^{-1}$ \\
$\eta$          & proportion of newborns female                & 0.5    & --- \\
$\gamma$        & progression rate $I\to A$ (Table value)      & 1.16   & yr$^{-1}$ \\
$\mu$           & natural (non-AIDS) death rate                & 0.0227 & yr$^{-1}$ \\
$\rho$          & social-mixing probability                    & 1.0    & --- \\
$\theta$        & female fecundity                             & 0.2088 & yr$^{-1}$ \\
$\vartheta$     & perinatal transmission probability           & 0.35   & --- \\
$\xi$           & ageing rate age-1$\to$age-2                   & 0.0667 & yr$^{-1}$ \\
$\zeta$         & proportion in sexual-activity class          & 1.0    & --- \\
\midrule
\multicolumn{4}{l}{\emph{Initial conditions (individuals):}
  $S_{f,1}=S_{m,1}=3000$, $S_{f,2}=S_{m,2}=1000$,} \\
\multicolumn{4}{l}{\quad all $I=0$, all $A=0$ except the seed (see
  \S\ref{sec:judgments}).} \\
\bottomrule
\end{tabular}
\end{table}

\paragraph{A conservation correction to the printed equations.}
The AIDS age-2 equation deserves a note. The textbook prints the ageing inflow to
$A_{f,2}$ (and $A_{m,2}$) as $\xi I_{f,1}$, but this violates conservation of
individuals. In Eq.~\ref{eq:Af1}, $A_{f,1}$ \emph{loses} $\xi A_{f,1}$ to ageing,
and that flux must land in $A_{f,2}$; meanwhile the ageing of $I_{f,1}$ is
\emph{already} counted as an inflow to $I_{f,2}$ in Eq.~\ref{eq:If2}. Using
$\xi I_{f,1}$ in Eq.~\ref{eq:Af2} would therefore double-count the $I_{f,1}$ cohort
and lose the $A_{f,1}$ cohort. I use the conservative form $\xi A_{f,1}$
(and $\xi A_{m,1}$) shown above. As a check, the total-population balance
$\frac{\mathrm{d}}{\mathrm{d}t}\sum(\text{states}) = \text{births}-\text{deaths}$
then holds to $\sim\!10^{-13}$ in the simulation, which would not be the case with
the printed term.

\subsection{Two modelling judgments}
\label{sec:judgments}

Two decisions were forced on me by trying to make the model reproduce the
textbook's stated behaviour. I present both honestly, because each changes the
results.

\paragraph{(i) The progression rate $\gamma$.}
For a two-sex frequency-dependent model the basic reproduction number is the
geometric mean of the two cross-sex reproduction numbers,
\begin{equation}
\Rzero = \frac{c\,\sqrt{\beta_{m,f}\,\beta_{f,m}}}{\mu+\gamma}.
\label{eq:R0}
\end{equation}
The Table~15.2 value $\gamma=1.16\,\text{yr}^{-1}$ implies a mean time from HIV
infection to AIDS of $1/(\mu+\gamma)\approx0.85$~yr --- less than a year, which
contradicts the textbook's own statement that progression takes 1--10~years. Worse,
substituting it into Eq.~\ref{eq:R0} gives $\Rzero\approx0.24<1$, so the epidemic
cannot ignite at all; a literal-parameter simulation confirms the seeded infection
dies out (peak prevalence $\sim\!7\times10^{-4}$). This directly contradicts
Fig.~15.5, where HIV persists endemically. I therefore adopt the biologically
grounded $\gamma=0.1\,\text{yr}^{-1}$, giving a mean infectious period of
$1/(\mu+\gamma)\approx8.1$~yr --- comfortably inside the textbook's stated 1--10~yr
range --- and
\begin{equation}
\Rzero = \frac{2.35\,\sqrt{0.2\times0.075}}{0.0227+0.1} \approx 2.35 > 1,
\end{equation}
so the epidemic invades, as the textbook figure requires. The associated
herd-immunity (critical vaccinated) fraction is $p_c = 1-1/\Rzero \approx 0.574$.

\paragraph{(ii) The epidemic seed.}
Table~15.2 seeds the epidemic with five AIDS males ($A_{m,2}=5$). But AIDS
individuals are excluded from the force of infection (Eq.~\ref{eq:foi}), so a seed
placed entirely in $A$ cannot transmit anything: the five seed males simply
progress to death without infecting anyone, and the epidemic never starts (verified
numerically --- maximum incidence exactly zero). I therefore seed five
\emph{infectious} males ($I_{m,2}=5$) instead. This is the smallest change that lets
the documented epidemic ignite while respecting the model's own rule that $A$ does
not transmit.

\subsection{Validation against Fig.~15.5}

With the two corrections above, the baseline ($\nu=0$) run reproduces the
qualitative behaviour of Fig.~15.5a (\verb|verify_baseline_prevalence.png|): the
total population grows over the early horizon (from $N(0)=8005$ to a peak of
$\sim\!22{,}900$ around year~48, driven by the high developing-country fertility),
HIV does \emph{not} burn out but settles to an endemic level, and female prevalence
exceeds male prevalence. The simulated equilibrium prevalences are
\begin{center}
female $\approx 0.77$, \quad male $\approx 0.61$, \quad overall $\approx 0.41$,
\end{center}
against the textbook targets of $\sim\!0.9$ (female) and $\sim\!0.75$ (male). The
ordering and endemic persistence match; the absolute levels run somewhat lower,
which is unsurprising given that I had to recalibrate $\gamma$ to make the epidemic
ignite at all. The time to exceed 10\% overall prevalence is about 34~yr, in the
same decades-long ballpark as the textbook's $\sim\!20$~yr.

\begin{figure}[t]
\centering
\includegraphics[width=0.78\linewidth]{verify_baseline_prevalence.png}
\caption{Baseline sIC validation ($\nu=0$). HIV persists endemically and female
prevalence exceeds male prevalence, reproducing the qualitative behaviour of
Fig.~15.5a; equilibrium prevalences are female $\approx0.77$, male $\approx0.61$.}
\label{fig:verify-baseline}
\end{figure}

The condom scenario (Fig.~15.5b) is modelled, as in the textbook, by halving both
transmission probabilities ($\beta_{f,m}:0.075\to0.0375$,
$\beta_{m,f}:0.2\to0.1$). This multiplies $\Rzero$ by exactly $0.5$ (it enters
Eq.~\ref{eq:R0} as $\sqrt{\beta_{m,f}\beta_{f,m}}$), giving $\Rzero\approx1.17$ ---
only just above one. In my calibration the resulting equilibrium prevalence is very
low ($\sim\!0.001$, \verb|verify_condom_prevalence.png|), which is a larger drop
than the textbook's Fig.~15.5b ($\sim\!0.7$ female, $\sim\!0.45$ male). The
direction is right (condoms lower prevalence and lengthen the time-to-10\%) but the
magnitude diverges, and the reason is instructive: at $\gamma=0.1$ the halved-$\beta$
$\Rzero$ sits barely above the threshold $\Rzero=1$, so the epidemic is in the
extremely sensitive near-threshold regime where small changes in calibration move
the endemic level a great deal. The textbook's larger residual prevalence is
consistent with its (implicitly) higher underlying $\Rzero$. I flag this divergence
rather than tune it away; it is the same near-threshold sensitivity that recurs in
Q(d).

\begin{figure}[t]
\centering
\includegraphics[width=0.78\linewidth]{verify_condom_prevalence.png}
\caption{Condom scenario (both $\beta$ halved). Prevalence collapses because
$\Rzero$ drops to $\approx1.17$, just above threshold. The direction matches
Fig.~15.5b; the magnitude is larger because the halved-$\beta$ system sits in the
sensitive near-threshold regime (see text).}
\label{fig:verify-condom}
\end{figure}

\subsection{The vaccination extension}

The vaccine adds two protected compartments, $P_{f,2}$ and $P_{m,2}$, at age~2 only
(age~2 is the terminal age class, so there is no ageing out of $P$). Susceptibles
are vaccinated at per-capita rate $\nu$ and protection wanes back to susceptible at
rate $l$:
\begin{align}
\dd{P_{f,2}}{t} &= \nu\,S_{f,2} - (l+\mu)\,P_{f,2}, &
\dd{P_{m,2}}{t} &= \nu\,S_{m,2} - (l+\mu)\,P_{m,2}, \\
\dd{S_{f,2}}{t} &= \xi S_{f,1} - (\lambda_{S_{f,2}}+\mu+\nu)\,S_{f,2} + l\,P_{f,2}, &
\dd{S_{m,2}}{t} &= \xi S_{m,1} - (\lambda_{S_{m,2}}+\mu+\nu)\,S_{m,2} + l\,P_{m,2},
\end{align}
with all other equations unchanged. A cost accumulator tracks the cumulative number
of vaccinations administered,
\begin{equation}
\dd{V}{t} = \nu\bigl(S_{f,2}+S_{m,2}\bigr),
\qquad \text{cost} = \$10\times V,
\label{eq:cost}
\end{equation}
integrated alongside the state so there is no post-hoc quadrature error. Each \$10
charge counts a vaccination \emph{event}, which correctly includes re-vaccinating
people whose protection waned and returned them to $S$. With $\nu=0$ the protected
and cost compartments stay identically zero and the model reduces exactly to the
baseline.

\paragraph{Vaccine type.}
I model a \emph{take} (all-or-nothing) vaccine with waning: while in $P$ an
individual is treated as fully protected (no force of infection acts on $P$, and
$P$ never flows to $I$), and a fraction loses protection at rate $l$. This is the
cleanest reading of the requested $S\to P\to S$ flow and maps onto the IC model's
``fully immunised + waning'' category. The project values $\nu=0.65\,\text{yr}^{-1}$
and $l=0.1\,\text{yr}^{-1}$ correspond to the Garnett/Stover ``65\% coverage'' and
``10-year mean duration'' standard scenario; using $\nu$ as a constant rate rather
than a coverage level reached after five years is a deliberate simplification.
Ignoring demography, the protected fraction saturates at a \emph{waning ceiling}
\begin{equation}
\frac{P}{S+P} \;\to\; \frac{\nu}{\nu+l} = \frac{0.65}{0.75} \approx 0.867,
\end{equation}
so even at this rate the vaccine can never protect everyone --- at most about 87\%
of the eligible pool at any instant.

\paragraph{The partner-pool convention.}
One modelling choice in the extension is genuinely load-bearing: whether the
protected individuals belong in the denominator of the frequency-dependent force of
infection. Protected people are sexually active and uninfected, so they \emph{are}
legitimate partners --- a partnership with a protected person simply transmits
nothing. The correct frequency-dependent denominator is therefore
$S+I+P$ (with $A$ still excluded). I adopt this convention. It matters because
moving susceptibles into $P$ then \emph{dilutes} the infectious frequency
$I/(S+I+P)$, which is exactly the mechanism of herd immunity. If instead one
excluded $P$ from the denominator, vaccinees would be removed from \emph{both} the
numerator and the denominator, artificially holding the infectious fraction up and
spuriously erasing the herd-immunity threshold --- vaccination would then appear to
have no critical rate at any level. That is an artifact of the denominator choice,
not a property of the disease; I report it only as a sensitivity check in Q(c). At
$\nu=0$ this choice is moot ($P=0$), so the baseline and condom scenarios are
unaffected.

\subsection{Q(a): Does vaccination cause a peak and decline?}

The question asks whether vaccination makes the epidemic peak and then decline. The
answer, at $\nu=0.65\,\text{yr}^{-1}$, is actually stronger than ``peak and
decline'': because the achievable protected fraction (0.867) comfortably exceeds the
herd-immunity threshold $p_c=0.574$, vaccination pushes the effective reproduction
number below one,
\begin{equation}
\Reff \approx \Rzero\Bigl(1-\tfrac{\nu}{\nu+l}\Bigr)
       \approx 2.35\times(1-0.867) \approx 0.31 < 1,
\end{equation}
so the seeded epidemic \emph{never ignites} in the first place.

The contrast with the baseline is stark (Table~\ref{tab:qa},
Figs.~\ref{fig:qa-incidence}--\ref{fig:qa-protected}). The baseline incidence rises
to a large peak of $\sim\!791$ new infections/yr around year~48 and stays high,
whereas under vaccination incidence is highest at $t=0$ (just the five seed males
infecting a handful of contacts) and then \emph{falls monotonically} to essentially
zero. Prevalence shows the same picture: the baseline climbs to a $\sim\!0.41$
endemic plateau, while the vaccinated run shows only a negligible interior bump
($\sim\!0.001$ around year~4.6) before declining to near zero. The decline is not
instantaneous because the existing $I$/$A$ stock carries the virus for the $\sim\!8$--$10$~yr
it takes to progress and die; but with $\Reff<1$ there is no fresh transmission to
refill it, so the infection clears. Meanwhile the protected fraction of the eligible
age-2 pool rises and saturates at $\approx0.85$, just below the analytic ceiling
$0.867$ (slightly lower because natural mortality $\mu$ also acts on $P$),
confirming the take-with-waning steady state.

\begin{table}[t]
\centering
\caption{Q(a): epidemic outcome with and without vaccination ($\nu=0.65$,
80-yr horizon).}
\label{tab:qa}
\begin{tabular}{lcc}
\toprule
 & Baseline ($\nu=0$) & Vaccine ($\nu=0.65$) \\
\midrule
Peak incidence (new/yr)      & 790.6 (@ 47.8 yr) & 2.3 (@ $t=0$, the seed) \\
Incidence behaviour          & rise to peak, stay high & monotone decline \\
Peak overall prevalence      & 0.409 & 0.0010 (@ $\approx$4.6 yr) \\
Prevalence @ 80 yr           & 0.408 & $\approx 0$ \\
Protected fraction $P/(S+P)$ & --- & $\to 0.85$ \\
\bottomrule
\end{tabular}
\end{table}

\begin{figure}[t]
\centering
\includegraphics[width=0.78\linewidth]{qa_incidence.png}
\caption{Q(a): incidence (new infections/yr), vaccine versus baseline. The baseline
peaks at $\sim\!791$/yr near year~48; under vaccination ($\Reff\approx0.31$)
incidence declines monotonically from the seed and never takes off.}
\label{fig:qa-incidence}
\end{figure}

\begin{figure}[t]
\centering
\includegraphics[width=0.78\linewidth]{qa_prevalence.png}
\caption{Q(a): overall prevalence, vaccine versus baseline. The baseline rises to a
$\sim\!0.41$ endemic plateau; the vaccinated run shows only a tiny early bump before
declining to near zero.}
\label{fig:qa-prevalence}
\end{figure}

\begin{figure}[t]
\centering
\includegraphics[width=0.78\linewidth]{qa_protected_fraction.png}
\caption{Q(a): protected fraction of the eligible age-2 pool, saturating near the
analytic waning ceiling $\nu/(\nu+l)=0.867$.}
\label{fig:qa-protected}
\end{figure}

\subsection{Q(b): Cost of vaccination}

Using the cost accumulator of Eq.~\ref{eq:cost}, Table~\ref{tab:qb} reports the
cumulative vaccinations, total cost, infections averted (cumulative baseline
infections minus cumulative vaccinated infections), and cost per infection averted
at several horizons.

\begin{table}[t]
\centering
\caption{Q(b): cost and cost-effectiveness of vaccination ($\nu=0.65$, \$10/vax).}
\label{tab:qb}
\begin{tabular}{rrrrr}
\toprule
$T$ (yr) & Vaccinations & Cost (\$) & Infect.\ averted & \$/infection averted \\
\midrule
20 & 11{,}374 & 113{,}738 &    305 & 372 \\
30 & 16{,}217 & 162{,}165 &  1{,}585 & 102 \\
50 & 24{,}746 & 247{,}461 & 13{,}434 &  18 \\
\bottomrule
\end{tabular}
\end{table}

Total cost grows nearly linearly: $\sim\!\$114$k at 20~yr, $\sim\!\$162$k at 30~yr,
$\sim\!\$247$k at 50~yr (Fig.~\ref{fig:qb}). Most of the expense is the up-front
vaccination of the standing susceptible pool; once the protected fraction
saturates, the programme keeps re-vaccinating waned individuals and new entrants at
a roughly steady $\sim\!407$ vaccinations/yr, i.e.\ a recurring cost of only
$\sim\!\$4{,}070$/yr.

Cost-effectiveness improves sharply over time. Early on the untreated baseline has
barely taken off, so there are few infections to avert and the ratio looks expensive
(\$372 per infection averted at 20~yr). As the baseline epidemic accelerates toward
its $\sim\!791$/yr incidence peak while the vaccinated arm stays near zero, averted
infections grow far faster than cost, and the cost per infection averted falls to
\$102 at 30~yr and \$18 at 50~yr. The 30-yr figure of \$102 sits just below the
\$110--390 band that Stover \& Garnett~(2002) report for the Imperial-College /
rural-Zimbabwe setting (on their \$20/person basis), which is consistent given the
\$10 unit cost here and the fact that, under the corrected partner-pool convention,
vaccination averts nearly the entire baseline epidemic rather than only part of it.
The absolute dollar totals depend on this small synthetic population, so the
transferable metric is the cost \emph{per infection averted}, not the headline cost.

\begin{figure}[t]
\centering
\includegraphics[width=0.78\linewidth]{qb_cost.png}
\caption{Q(b): (top) cumulative cost and vaccinations versus time, with horizon
markers at 20/30/50~yr; (bottom) cumulative infections for baseline versus vaccine
and the infections-averted band.}
\label{fig:qb}
\end{figure}

\subsection{Q(c): Optimum vaccination rate}

``Optimum'' has two defensible meanings, and I report both, because they give
different answers.

\paragraph{Epidemiological threshold $\nu_c$.}
The smallest rate that drives $\Reff<1$ follows from requiring the protected
fraction to reach $p_c$. Since the steady protected fraction is $\nu/(\nu+l)$, the
condition $\nu/(\nu+l)\ge p_c$ gives a critical rate
\begin{equation}
\nu_c = \frac{l\,p_c}{1-p_c},
\qquad p_c = 1-\tfrac{1}{\Rzero}\approx0.574,
\end{equation}
which evaluates to $\nu_c\approx0.135\,\text{yr}^{-1}$. This closed form is the
right order of magnitude and direction but is a lower bound, because it omits two
mechanisms the full sIC model contains: perinatal transmission ($\vartheta=0.35$),
an infection route vaccination does not block, and demographic recruitment with
mortality acting on $P$ (which holds the realised protected fraction slightly below
the ceiling). Measuring the threshold directly --- seeding a tiny infection into the
vaccinated disease-free population and finding the $\nu$ at which the early growth
rate crosses zero ($\Reff=1$) --- gives a simulated $\nu_c\approx0.366\,\text{yr}^{-1}$
(Fig.~\ref{fig:qc-threshold}). The standard $\nu=0.65$ sits comfortably above this,
so it genuinely controls the epidemic. I emphasise that this is a frequency-dependent
\emph{prevalence-proportion} threshold rather than an ``$(I+A)/N\to0$'' criterion:
the severe AIDS mortality ($\alpha=1\,\text{yr}^{-1}$) drives the whole population
toward extinction over centuries even at $\nu=0$, so absolute counts collapse for
the wrong reason, and the invasion / $\Reff$ test is the robust signal.

\paragraph{Cost-effective rate.}
Sweeping $\nu$ over a 50-yr horizon, infections averted \emph{saturate} past
$\nu_c$: below it each extra unit of $\nu$ averts many infections, but once
$\nu\gtrsim\nu_c$ essentially all of the achievable $\sim\!13{,}400$ infections are
already averted and further vaccination only re-vaccinates waned individuals for
negligible benefit (Figs.~\ref{fig:qc-prev},~\ref{fig:qc-ce}). The
diminishing-returns knee --- the smallest $\nu$ capturing 99\% of the maximum
achievable aversion --- is at $\nu\approx0.175\,\text{yr}^{-1}$, with a marginal cost
there of $\sim\!\$122$ per extra infection averted. The standard $\nu=0.65$ still
controls the epidemic (it is above $\nu_c$) but over-vaccinates somewhat relative to
this cost-effective knee; its average cost is \$18 per infection averted at 50~yr.

\paragraph{Sensitivity.}
The exact threshold tracks $\Rzero$: raising the partner rate to $c=3.0$ lifts
$\Rzero$ to $\sim\!3.0$ and the simulated $\nu_c$ to $\sim\!0.48$; lowering it via
$\gamma=0.18$ drops $\Rzero$ to $\sim\!1.42$ and $\nu_c$ to $\sim\!0.16$. The
conclusion is robust in direction --- a finite $\nu_c=l\,p_c/(1-p_c)$ always exists
under the correct convention and $\nu=0.65$ stays above it across this plausible
range --- but near-threshold cases (small $\Rzero$) leave the margin thin.

\paragraph{Caveat.}
If $P$ is wrongly excluded from the partner pool, the invasion growth rate never
crosses zero in the sweep: vaccination would appear to have \emph{no} finite
threshold at any rate and merely shrink the epidemic. As argued above this is an
artifact of dropping uninfected, sexually-active vaccinees out of the
frequency-dependent denominator, and I report it only to show that the threshold
result is genuine, not an accident of the denominator choice.

\begin{figure}[t]
\centering
\includegraphics[width=0.80\linewidth]{qc_invasion_threshold.png}
\caption{Q(c): invasion growth rate $r(\nu)$ crossing zero ($\Reff=1$) under the
corrected convention ($P$ in pool, finite $\nu_c\approx0.366$); the wrong
$P$-excluded convention (no crossing) is shown for contrast. Theory $\nu_c$,
simulated $\nu_c$, and $\nu=0.65$ are marked.}
\label{fig:qc-threshold}
\end{figure}

\begin{figure}[t]
\centering
\includegraphics[width=0.78\linewidth]{qc_prevalence_vs_nu.png}
\caption{Q(c): peak epidemic size versus $\nu$. Above the threshold the epidemic is
suppressed.}
\label{fig:qc-prev}
\end{figure}

\begin{figure}[t]
\centering
\includegraphics[width=0.78\linewidth]{qc_cost_effectiveness.png}
\caption{Q(c): average and marginal (incremental) cost per infection averted versus
$\nu$, with the diminishing-returns knee at $\nu\approx0.175$ marked.}
\label{fig:qc-ce}
\end{figure}

\subsection{Q(d): Vaccination versus condoms}

Finally, which intervention controls the epidemic better: vaccination ($\nu=0.65$,
$l=0.1$) or safe-sex education and condoms (modelled, as in the textbook, by halving
both transmission probabilities)? They act on different levers --- condoms lower
$\Rzero$ directly through $\beta$, vaccination moves people into $P$ --- so I compare
them on matched outcomes (Table~\ref{tab:qd}, Figs.~\ref{fig:qd-comp},~\ref{fig:qd-averted}).

\begin{table}[t]
\centering
\caption{Q(d): strategy comparison (80-yr horizon, $\Rzero=2.35$ baseline).}
\label{tab:qd}
\begin{tabular}{lcrr}
\toprule
Strategy & $\Rzero$ / $\Reff$ & Infect.\ averted & Cost (\$) \\
\midrule
Baseline                       & $\Rzero=2.35$         & ---       & 0 \\
Vaccination ($\nu=0.65$)       & $\Reff=0.31$          & 31{,}224  & 348{,}618 \\
Condoms ($0.5\beta$)           & $\Rzero=1.17$         & 30{,}857  & 0 (unpriced) \\
\bottomrule
\end{tabular}
\end{table}

Both strategies essentially control the epidemic. Vaccination drives
$\Reff=0.31<1$ (invasion rate $-0.040\,\text{yr}^{-1}$): it strictly crosses the
elimination threshold because the protected fraction $0.867$ exceeds $p_c=0.574$.
Condoms put $\Rzero$ at $1.17$ --- still marginally \emph{above} one (invasion rate
$+0.086\,\text{yr}^{-1}$), so a 50\% $\beta$ cut is not quite enough to cross the
threshold at this $\Rzero$, though the epidemic is so slow that it is effectively
controlled within the horizon. On infections averted the two are nearly identical:
vaccination averts $31{,}224$ of the baseline $31{,}235$, condoms $30{,}857$, with
vaccination edging ahead only because it pushes $\Reff$ strictly below one (residual
$11$ cumulative infections versus $378$ for condoms). Both collapse peak prevalence
from $0.41$ to $\sim\!0.001$ with no large interior peak.

Where they differ sharply is cost basis, and this is where I decline to crown a
single winner. Vaccination carries an \emph{explicit} \$348{,}618 over 80~yr and is
capped by the waning ceiling $0.867$, so it must re-vaccinate waned individuals
indefinitely. Condoms are modelled as a permanent $0.5\times$ cut to $\Rzero$ at no
priced cost --- but that cost is \emph{unpriced}, not zero; a real behaviour-change
programme has its own budget the project does not supply. A dollar-matched
comparison would need a condom-programme cost we do not have. There is also the
near-threshold caveat: condoms leave $\Rzero$ only $\sim\!17\%$ above one, and a
higher baseline $\Rzero$ (say a higher partner rate) would push condom-$\Rzero$
clearly above one while raising vaccination's $\nu_c$ --- both verdicts are
$\Rzero$-sensitive near threshold, which is exactly why I report $\Reff$ explicitly
rather than just trajectories. The honest reading is a near-tie on epidemiological
outcome, with vaccination marginally ahead for strictly crossing the threshold, and
the two not directly comparable on cost as priced here.

\begin{figure}[t]
\centering
\includegraphics[width=0.78\linewidth]{qd_strategy_comparison.png}
\caption{Q(d): overall prevalence trajectories for baseline, vaccination, and
condoms. Both interventions collapse the epidemic; the baseline rises to a
$\sim\!0.41$ plateau.}
\label{fig:qd-comp}
\end{figure}

\begin{figure}[t]
\centering
\includegraphics[width=0.78\linewidth]{qd_averted_and_reff.png}
\caption{Q(d): infections averted and reproduction number per strategy. Vaccination
crosses $\Reff<1$ strictly; condoms sit just above $\Rzero=1$.}
\label{fig:qd-averted}
\end{figure}

\subsection{Part 2 summary}

On the sIC model, with $\gamma$ recalibrated to a biologically sensible
$0.1\,\text{yr}^{-1}$ (giving $\Rzero\approx2.35$) and the epidemic seeded in the
infectious rather than the AIDS compartment, vaccination at $\nu=0.65\,\text{yr}^{-1}$
controls the HIV epidemic: it drives $\Reff\approx0.31<1$, so the epidemic never
ignites, and prevalence stays near zero instead of rising to the $\sim\!0.41$
endemic plateau of the baseline. It is cost-effective ($\sim\!\$102$ per infection
averted at 30~yr, within the published Garnett/Stover range), though the standard
rate over-vaccinates somewhat relative to the cost-effective knee at
$\nu\approx0.18\,\text{yr}^{-1}$. Against condom promotion the two are roughly
comparable on epidemiological outcome --- both essentially control the epidemic,
vaccination crossing the threshold strictly while condoms hover just above it --- but
they are not directly comparable on cost as priced here, and both verdicts are
sensitive to the baseline $\Rzero$ near threshold. All of these conclusions rest on
the documented modelling assumptions, in particular the $\gamma$ recalibration and
the partner-pool convention that gives vaccination its herd-immunity effect.

\clearpage

\section{Reproducibility and code}
All models and analyses are in Python (NumPy/SciPy/Matplotlib). Part~1 lives in
\texttt{part1\_glucose/} (\texttt{cobelli.py} core; \texttt{ex1\_ivgtt.py}\,$\dots$\,\texttt{ex7\_icecream.py};
\texttt{verify\_core.py}); Part~2 in \texttt{part2\_hiv/} (\texttt{sic.py} core; \texttt{q\_a}\,$\dots$\,\texttt{q\_d}
analyses; \texttt{verify\_sic.py}). Each core ships a verification harness that checks the basal
steady state / conservation of individuals and agreement between the LSODA and Radau integrators.
Octave (\texttt{.m}) cross-check scripts for the two cores are provided in \texttt{octave\_crosscheck/}.

\section{References}
\small
\begin{enumerate}[leftmargin=1.4em,itemsep=1pt]
\item Cobelli, C., Federspil, G., Pacini, G., Salvan, A., \& Scandellari, C. (1982). An integrated
mathematical model of the dynamics of blood glucose and its hormonal control. \emph{Mathematical
Biosciences}, 58, 27--60.
\item Makroglou, A., Li, J., \& Kuang, Y. (2006). Mathematical models and software tools for the
glucose--insulin regulatory system and diabetes: an overview. \emph{Applied Numerical Mathematics},
56(3--4), 559--573.
\item Dalla Man, C., Rizza, R.~A., \& Cobelli, C. (2007). Meal simulation model of the
glucose--insulin system. \emph{IEEE Transactions on Biomedical Engineering}, 54(10), 1740--1749.
\item Garnett, G.~P., \& Anderson, R.~M. (1993). Factors controlling the spread of HIV in
heterosexual communities in developing countries. \emph{Philosophical Transactions of the Royal
Society of London B}, 342, 137--159.
\item Garnett, G.~P., Desai, K., \& Williams, J. (2002). The epidemiological impact of an HIV/AIDS
vaccination as a function of vaccine properties (Imperial College model). In J.~Stover \emph{et~al.}
(Eds.), \emph{Policy Research Working Paper 2811}. World Bank.
\item Haefner, J.~W. (2005). \emph{Modeling Biological Systems: Principles and Applications}
(2nd~ed., Ch.~12 \& 15). Springer.
\end{enumerate}

\end{document}
